% =====================================================================
%  Projeto FEMM – Motores de Corrente Alternada  (versão corrigida)
%  Q_s=36, Q_r=28, q=2
%  Disciplina: Conversão da Energia – UFMG 2026/1
%  Prof. Braz J. Cardoso Filho, Ph.D.
%  Autores: Bruno Ribeiro e Arnaldo
% =====================================================================
\documentclass[12pt,a4paper]{article}

\usepackage[utf8]{inputenc}
\usepackage[T1]{fontenc}
\usepackage[brazil]{babel}
\usepackage{mathptmx}            % Times New Roman (texto e matemática)
\usepackage[a4paper,
            top=2.0cm, bottom=2.0cm,
            left=2.0cm, right=2.0cm,
            headheight=14pt]{geometry}
\usepackage{tabularx}

% ── Parágrafos: 3 pt de separação, sem recuo
\usepackage{parskip}
\setlength{\parskip}{3pt}
\setlength{\parindent}{0pt}

% ── Cabeçalho
\usepackage{fancyhdr}
\pagestyle{fancy}
\fancyhf{}
\renewcommand{\headrulewidth}{0pt}
\fancyhead[L]{\small Projeto FEMM -- Análise Harmônica de Motor CA -- UFMG}
\fancyhead[R]{\small B. Ribeiro \& Arnaldo\quad\thepage}

\usepackage{amsmath,amssymb}
\usepackage{booktabs}
\usepackage{multirow}
\usepackage{graphicx}
\usepackage{caption}
\usepackage{siunitx}
\usepackage{titlesec}
\usepackage{enumitem}

\titleformat{\section}{\normalsize\bfseries}{\thesection.}{0.4em}{}
\titlespacing*{\section}{0pt}{5pt}{2pt}
\titleformat{\subsection}{\normalsize\itshape}{\thesubsection.}{0.4em}{}
\titlespacing*{\subsection}{0pt}{3pt}{1pt}

\captionsetup{font=small, labelsep=period, skip=2pt}
\sisetup{output-decimal-marker={,}}
\renewcommand{\arraystretch}{1.07}

% =====================================================================
\begin{document}
\fontsize{12}{15}\selectfont
% =====================================================================

{\centering\bfseries
Análise harmônica de uma máquina de indução trifásica
via elementos finitos (FEMM)\\[2pt]}
{\centering\small
Bruno Ribeiro \& Arnaldo \;|\;
Engenharia Elétrica -- UFMG \;|\;
Prof.\ Braz J.\ Cardoso Filho -- Conversão da Energia -- 2026/1\par}
\vspace{3pt}

% =====================================================================
\section{Dados do projeto}
% =====================================================================

\begin{table}[h!]
\vspace{-2pt}\centering\small
\caption{Parâmetros da máquina (extraídos do enunciado e da laminação Fig.~1).}
\label{tab:dados}
\setlength{\tabcolsep}{5pt}
\begin{tabular}{llll}
\toprule
\textbf{Parâmetro} & \textbf{Valor} & \textbf{Parâmetro} & \textbf{Valor}\\
\midrule
Polos $P$ / pares $p$   & $6$ / $3$    & Frequência $f_e$       & \SI{60}{Hz}\\
Fases $m$               & $3$          & Tensão terminal $V$    & \SI{127}{V_{rms}}\\
Ranhuras estator $Q_s$  & $36$         & Entreferro $g$         & \SI{0{,}5}{mm}\\
Ranhuras rotor $Q_r$    & $28$         & Comprimento pacote $L$ & \SI{140}{mm}\\
Ranhuras/polo/fase $q$  & $2$          & Aço laminações         & M250-50A\\
Passo polar $\tau$      & $6$ ranhuras & $B_{\max}$ no núcleo   & \SI{1{,}7}{T}\\
$\alpha$ (el./ranhura)  & \SI{30}{\degree} & $B_{g,1}$ alvo     & \SI{0{,}9}{T}\\
Diâm.\ interno estator  & \SI{115}{mm} & Velocidade síncrona    & \SI{1200}{rpm}\\
\bottomrule
\end{tabular}
\end{table}
\vspace{-2pt}

Analisam-se três configurações de enrolamento com ligação em série e
1~cinturão por polo por fase ($q{=}2$ ranhuras/polo/fase):
\textbf{(I)}~camada simples, passo pleno ($y_1{=}6$, $\beta{=}1$,
\SI{180}{\degree}~el.) --- 2~bobinas por polo por fase;
\textbf{(II)}~camada dupla, passo pleno (mesma distribuição eletromagnética
de~I, mas com 4~bobinas por polo por fase pela segunda camada);
\textbf{(III)}~camada dupla, passo encurtado $\beta{=}5/6$
($y_1{=}5$ ranhuras, \SI{150}{\degree}~el.), também com 4~bobinas por polo
por fase --- escolha clássica para atenuar a quinta e a sétima harmônica.

O número de espiras em série por fase é $N_f{=}(Q_s/2m)\,N_c$
para camada simples e $N_f{=}(Q_s/m)\,N_c$ para camada dupla,
onde $N_c{=}10$ é o número de condutores por ranhura por camada:

\begin{equation}
N_f^{\,(I)} = \frac{36}{2\times3}\times10 = 60\;\text{espiras},
\qquad
N_f^{\,(II,III)} = \frac{36}{3}\times10 = 120\;\text{espiras}.
\label{eq:Nf}
\end{equation}

% =====================================================================
\section{Fatores $k_p$, $k_d$ e $k_w$ -- item (a)}
% =====================================================================

Para enrolamento distribuído em $q$ ranhuras equiespaçadas por polo por fase,
com ângulo elétrico $\alpha$ entre ranhuras adjacentes e passo relativo
$\beta{=}y_1/\tau$, os fatores de distribuição, passo e enrolamento para o
harmônico de ordem~$\nu$ são:
%
\begin{equation}
k_{d,\nu}=\frac{\sin\!\bigl(\nu\,q\,\alpha/2\bigr)}{q\,\sin\!\bigl(\nu\,\alpha/2\bigr)},
\qquad
k_{p,\nu}=\sin\!\left(\nu\beta\,\frac{\pi}{2}\right),
\qquad
k_{w,\nu}=k_{p,\nu}\cdot k_{d,\nu}.
\label{eq:kw}
\end{equation}

Com $Q_s{=}36$, $P{=}6$, $m{=}3$: $\alpha{=}P\cdot180^{\circ}/Q_s{=}30^{\circ}$
e $q{=}Q_s/(Pm){=}2$.

\subsection*{Cálculo numérico para $\nu=1$}

Substituindo $q{=}2$ e $\alpha{=}30^{\circ}$ em~\eqref{eq:kw}:
\[
k_{d,1}=\frac{\sin(2\times30^{\circ}/2)}{2\sin(30^{\circ}/2)}
        =\frac{\sin 30^{\circ}}{2\sin 15^{\circ}}
        =\frac{0{,}5000}{2\times0{,}2588}
        =0{,}9659.
\]
Para passo pleno ($\beta{=}1$): $k_{p,1}=\sin 90^{\circ}=1{,}0000$,
logo $k_{w,1}^{\,\beta=1}=0{,}9659$.
Para passo encurtado ($\beta{=}5/6$):
$k_{p,1}=\sin\!\bigl(\tfrac{5}{6}{\times}90^{\circ}\bigr)=\sin 75^{\circ}=0{,}9659$,
logo $k_{w,1}^{\,\beta=5/6}=0{,}9659\times0{,}9659=0{,}9330$
(redução de apenas $3{,}4\%$ na fundamental).

Para $\nu{=}5$ e $\nu{=}7$ com $\beta{=}5/6$:
$k_{p,5}=\sin(5{\times}75^{\circ})=\sin 375^{\circ}=\sin 15^{\circ}\approx0{,}259$
e $k_{p,7}=\sin(7{\times}75^{\circ})=\sin 525^{\circ}=\sin 165^{\circ}\approx0{,}259$.
O fator de passo reduz ambos para apenas $25{,}9\%$ do valor de passo pleno,
resultando em $|k_{w,5}^{\,5/6}|{=}0{,}259\times0{,}259{=}0{,}067$
($-74\%$) e $|k_{w,7}^{\,5/6}|{=}0{,}067$ ($-74\%$).

\begin{table}[h!]
\vspace{-1pt}\centering\small
\caption{Fatores de enrolamento para $q{=}2$, $\alpha{=}30^{\circ}$.
         Triplens omitidas (canceladas pela soma trifásica equilibrada).}
\label{tab:fatores}
\setlength{\tabcolsep}{5pt}
\begin{tabular}{crrrrrr}
\toprule
$\nu$ & Sentido & $|k_{d,\nu}|$ & $k_{p,\nu}\ (\beta{=}1)$
      & $k_{p,\nu}\ (\beta{=}5/6)$
      & $|k_{w,\nu}|\ (\beta{=}1)$ & $|k_{w,\nu}|\ (\beta{=}5/6)$\\
\midrule
 1 & direto  & $0{,}9659$ & $1{,}0000$ & $0{,}9659$ & $0{,}9659$ & $0{,}9330$\\
 5 & reverso & $0{,}2588$ & $1{,}0000$ & $0{,}2588$ & $0{,}2588$ & $0{,}0670$\\
 7 & direto  & $0{,}2588$ & $1{,}0000$ & $0{,}2588$ & $0{,}2588$ & $0{,}0670$\\
11 & reverso & $0{,}9659$ & $1{,}0000$ & $0{,}9659$ & $0{,}9659$ & $0{,}9330$\\
13 & direto  & $0{,}9659$ & $1{,}0000$ & $0{,}9659$ & $0{,}9659$ & $0{,}9330$\\
\bottomrule
\end{tabular}
\end{table}

Uma consequência importante de $q{=}2$ é que o fator de distribuição assume
apenas dois valores: $|k_{d,\nu}|{=}0{,}966$ para $\nu\equiv\pm1\pmod{12}$
e $|k_{d,\nu}|{=}0{,}259$ para $\nu\equiv\pm5\pmod{12}$.
Em particular, $|k_{d,11}|{=}|k_{d,13}|{=}|k_{d,1}|{=}0{,}966$ ---
diferentemente do caso $q{=}4$ ($Q_s{=}72$), onde esses harmônicos eram
fortemente atenuados ($|k_{d,11}|{=}0{,}126$).
Assim, $\nu{=}11$ e $\nu{=}13$ \emph{não} são suprimidos pelo enrolamento.
O encurtamento $5/6$ também não os atenua: $|k_{p,11}|{=}|k_{p,13}|{=}0{,}966$.

% =====================================================================
\section{Espectro analítico de $B_g(\theta)$ -- item (b)}
% =====================================================================

\subsection*{Corrente de pico e fórmula do campo}

Assumindo $\mu_\text{ferro}\to\infty$ e entreferro uniforme, a amplitude da
$\nu$-ésima harmônica espacial da densidade de fluxo no entreferro de uma
máquina trifásica excitada em $t{=}0$ com $i_A{=}I_\text{pk}$,
$i_B{=}i_C{=}{-}I_\text{pk}/2$ é (Fitzgerald, eq.~4.6):
%
\begin{equation}
|B_{g,\nu}| = \frac{6\mu_0}{\pi}\,
              \frac{k_{w,\nu}\,N_f\,I_\text{pk}}{P\,g\,\nu},
\label{eq:Bgnu}
\end{equation}
%
donde, impondo $|B_{g,1}|{=}B_{g,1}^{\star}{=}\SI{0{,}9}{T}$:
%
\begin{equation}
N_f\,I_\text{pk} = \frac{B_{g,1}^{\star}\,\pi\,P\,g}{6\,\mu_0\,k_{w,1}},
\qquad
\frac{|B_{g,\nu}|}{|B_{g,1}|}=\frac{|k_{w,\nu}|}{\nu\,|k_{w,1}|}.
\label{eq:espectro}
\end{equation}

As correntes de pico para cada configuração são:

\begin{equation}
I_\text{pk}^{(I)}
  = \frac{B_{g,1}^{\star}\,\pi\,P\,g}{6\,\mu_0\,k_{w,1}^{\beta=1}\,N_f^{(I)}}
  = \frac{0{,}9\times\pi\times6\times5\times10^{-4}}
         {6\times4\pi\times10^{-7}\times0{,}9659\times60}
  \approx \SI{19{,}4}{A},
\label{eq:Ipk1}
\end{equation}
\[
I_\text{pk}^{(II)} = \frac{N_f^{(I)}\,I_\text{pk}^{(I)}}{N_f^{(II)}}
  = \frac{60\times19{,}4}{120} \approx \SI{9{,}70}{A},
\qquad
I_\text{pk}^{(III)} \approx \SI{10{,}05}{A}.
\]

\begin{table}[h!]
\vspace{-1pt}\centering\small
\caption{Espectro analítico de $B_g$ com $|B_{g,1}|{=}\SI{0{,}9}{T}$
         ($\mu_\text{ferro}\to\infty$, entreferro liso).
         Triplens canceladas. Harmônicos de ranhura \emph{não} modelados aqui.}
\label{tab:espectro}
\setlength{\tabcolsep}{6pt}
\begin{tabular}{rccccc}
\toprule
      & \multicolumn{2}{c}{$\beta{=}1$ (Configs~I e~II)}
      & \multicolumn{2}{c}{$\beta{=}5/6$ (Config~III)}\\
\cmidrule(lr){2-3}\cmidrule(lr){4-5}
$\nu$ & $|k_{w,\nu}|/(\nu|k_{w,1}|)$ & $|B_{g,\nu}|$\,(T)
      & $|k_{w,\nu}|/(\nu|k_{w,1}|)$ & $|B_{g,\nu}|$\,(T)\\
\midrule
 1 & $1{,}0000$ & $0{,}9000$ & $1{,}0000$ & $0{,}9000$\\
 5 & $0{,}0536$ & $0{,}0482$ & $0{,}0144$ & $0{,}0129$\\
 7 & $0{,}0383$ & $0{,}0344$ & $0{,}0103$ & $0{,}0092$\\
11 & $0{,}0909$ & $0{,}0818$ & $0{,}0909$ & $0{,}0818$\\
13 & $0{,}0769$ & $0{,}0692$ & $0{,}0769$ & $0{,}0692$\\
\bottomrule
\end{tabular}
\end{table}

O encurtamento $5/6$ reduz $|B_{g,5}|$ de \SI{48{,}2}{mT} para \SI{12{,}9}{mT}
($-73\%$) e $|B_{g,7}|$ de \SI{34{,}4}{mT} para \SI{9{,}2}{mT} ($-73\%$),
mantendo a quinta e sétima harmônicas abaixo de $1{,}5\%$ de $B_{g,1}$.
Em contrapartida, $|B_{g,11}|$ e $|B_{g,13}|$ permanecem em $9{,}1\%$ e $7{,}7\%$
de $B_{g,1}$ em todas as configurações --- reflexo direto do menor espalhamento
do enrolamento ($q{=}2$) em relação ao projeto com $q{=}4$.

Adicionalmente, com $Q_s{=}36$ e $p{=}3$, os harmônicos de ranhura do estator
concentram-se em ordens $\nu{=}k(Q_s/p){\pm}1{=}12k{\pm}1$, ou seja
$\nu{=}11,13,23,25,\ldots$ As ordens $\nu{=}11$ e $\nu{=}13$ coincidem
simultaneamente com harmônicos de enrolamento e harmônicos de ranhura;
na simulação por elementos finitos (itens c e~d) espera-se que as amplitudes
nessas ordens superem o analítico por receberem contribuição de ambas as fontes.

% =====================================================================
\begin{thebibliography}{9}\small\vspace{-2pt}
\bibitem{fitz}
  A.~E.~Fitzgerald, C.~Kingsley~Jr., S.~D.~Umans,
  \textit{Máquinas Elétricas}, 7\textsuperscript{a}~ed.\ AMGH, 2014.
\bibitem{chap}
  S.~J.~Chapman, \textit{Fundamentos de Máquinas Elétricas},
  5\textsuperscript{a}~ed.\ AMGH, 2013.
\bibitem{femm}
  D.~C.~Meeker,
  \textit{Finite Element Method Magnetics -- User's Manual}, v.~4.2, 2019.
  Disponível em: \texttt{femm.info}.
\end{thebibliography}

\end{document}
